% PRIVATE -- os/samir/cmd/proc.c p_private -> xb_interp_declare_private
% (hides the caller's variable until the routine returns).  Gate:
% test-interp-proc.
\begin{XbaseEntry}{PRIVATE}{command}
Creates memory variables that belong to the procedure running, hiding any of
the same names.
\Syntax
\SyntaxLine{\kw{PRIVATE} \xp{memvar list}}
\Usage
\cmd{PRIVATE} gives the procedure in which it is used its own memory variables
of the names in the list.  A memory variable of the same name created by the
program that called the procedure, or declared PUBLIC, is hidden while the
procedure runs and keeps its value: when the procedure returns, the private
variables are released and the hidden ones are visible again.

A private memory variable has no value until a value is stored in it.

A memory variable that a procedure creates without \cmd{PRIVATE}, by storing a
value in a name that does not yet exist, is private to that procedure as well.
\Examples
\file{TALLY.PRG} keeps its count apart from the count of its procedure:
\ProgramListing{TALLY.PRG}
\DotExample{examples/c-private}
\SeeAlso{PUBLIC, RELEASE, STORE, PARAMETERS}
\end{XbaseEntry}

